\documentclass[aspectratio=169]{beamer}
\usepackage[english]{babel}
\usepackage[utf8]{inputenc}
\usepackage[T1]{fontenc}
\usepackage{lmodern}
\usepackage{listings}
\usepackage{xcolor}
\usepackage{graphicx}
\usepackage{textcomp}
\DeclareUnicodeCharacter{2014}{---}
\DeclareUnicodeCharacter{2013}{--}
\DeclareUnicodeCharacter{2212}{-}
\DeclareUnicodeCharacter{2192}{\ensuremath{\rightarrow}}
\DeclareUnicodeCharacter{00B1}{\ensuremath{\pm}}
\DeclareUnicodeCharacter{00D7}{\ensuremath{\times}}
\DeclareUnicodeCharacter{00B7}{\ensuremath{\cdot}}
\DeclareUnicodeCharacter{20AC}{EUR}
\DeclareUnicodeCharacter{2070}{\textsuperscript{0}}
\DeclareUnicodeCharacter{00B9}{\textsuperscript{1}}
\DeclareUnicodeCharacter{00B2}{\textsuperscript{2}}
\DeclareUnicodeCharacter{00B3}{\textsuperscript{3}}
\DeclareUnicodeCharacter{2074}{\textsuperscript{4}}
\DeclareUnicodeCharacter{2075}{\textsuperscript{5}}
\DeclareUnicodeCharacter{2076}{\textsuperscript{6}}
\DeclareUnicodeCharacter{2077}{\textsuperscript{7}}
\DeclareUnicodeCharacter{2078}{\textsuperscript{8}}
\DeclareUnicodeCharacter{2079}{\textsuperscript{9}}
\DeclareUnicodeCharacter{2248}{\ensuremath{\approx}}
\DeclareUnicodeCharacter{2264}{\ensuremath{\le}}

% Colors for code
\definecolor{codebg}{RGB}{245,245,245}
\definecolor{codecomment}{RGB}{0,128,0}
\definecolor{codekeyword}{RGB}{0,0,255}
\definecolor{codestring}{RGB}{163,21,21}
\definecolor{bandcolor}{RGB}{200,50,50}

\lstdefinestyle{pythonstyle}{
    language=Python,
    backgroundcolor=\color{codebg},
    basicstyle=\ttfamily\scriptsize,
    keywordstyle=\color{codekeyword},
    commentstyle=\color{codecomment},
    stringstyle=\color{codestring},
    showstringspaces=false,
    breaklines=true,
    frame=single,
    rulecolor=\color{black},
    escapeinside={(*@}{@*)},
    moredelim=**[l][\color{bandcolor}\bfseries]{\#\ ---\ NEW},
    moredelim=**[l][\color{bandcolor}\bfseries]{\#\ ---\ CHANGED},
    literate=
      {á}{{\'a}}1 {é}{{\'e}}1 {í}{{\'i}}1 {ó}{{\'o}}1 {ú}{{\'u}}1
      {Á}{{\'A}}1 {É}{{\'E}}1 {Í}{{\'I}}1 {Ó}{{\'O}}1 {Ú}{{\'U}}1
      {ñ}{{\~n}}1 {Ñ}{{\~N}}1 {ü}{{\"u}}1 {Ü}{{\"U}}1
      {¿}{{?`}}1 {¡}{{!`}}1
      {—}{{---}}1 {–}{{--}}1 {−}{{-}}1
      {±}{{\ensuremath{\pm}}}1 {×}{{\ensuremath{\times}}}1
      {→}{{\ensuremath{\rightarrow}}}1
      {°}{{\textdegree}}1
      {·}{{\ensuremath{\cdot}}}1
      {€}{{\texteuro}}1
      {≈}{{\ensuremath{\approx}}}1
      {≤}{{\ensuremath{\le}}}1
      {⁰}{{\textsuperscript{0}}}1 {¹}{{\textsuperscript{1}}}1
      {²}{{\textsuperscript{2}}}1 {³}{{\textsuperscript{3}}}1
      {⁴}{{\textsuperscript{4}}}1 {⁵}{{\textsuperscript{5}}}1
      {⁶}{{\textsuperscript{6}}}1 {⁷}{{\textsuperscript{7}}}1
      {⁸}{{\textsuperscript{8}}}1 {⁹}{{\textsuperscript{9}}}1
}

\lstset{style=pythonstyle}

% Title slide macro: \lessontitle{Lesson Number}{Title}
\newcommand{\lessontitle}[2]{
    \title{Lesson #1: #2}
    \author{}
    \institute{Tec de Monterrey}
    \begin{frame}[plain]
        \titlepage
    \end{frame}
}

% Figure frame macro: \figureframe{Title}{Image path}
\newcommand{\figureframe}[2]{
    \begin{frame}{#1}
        \begin{center}
            \includegraphics[height=0.8\textheight,width=\textwidth,keepaspectratio]{#2}
        \end{center}
    \end{frame}
}

\setbeamertemplate{navigation symbols}{}
\setbeamertemplate{footline}[frame number]
\date{}
\begin{document}
\lessontitle{05}{Mutation}
\begin{frame}{Design question}
\normalsize
How far does one mutation move a chromosome? Rate and step size are separate choices.
\end{frame}
\begin{frame}{Real-valued model}
\normalsize
For each gene, \(B_i\sim\mathrm{Bernoulli}(p)\), \(\Delta_i\sim\mathcal N(0,\sigma^2)\), \(x_i'=\Pi_{[l,u]}(x_i+B_i\Delta_i)\). Expected activated genes: \(Lp\).
\end{frame}
\begin{frame}{Rate and step}
\normalsize
With \(L=48\) and \(p=.10\), 4.8 genes activate on average. Holding \(p\sigma\) fixed does not fix the longest jump or the chance to cross a low-fitness region.
\end{frame}
\begin{frame}{Three regimes}
\normalsize
\begin{center}\includegraphics[width=.84\textwidth,height=.53\textheight,keepaspectratio]{../../figures/mutation_reach_02_rate_and_step.png}\end{center}
{\small Compare touched genes, total distance, and longest jump, not one output alone.}
\end{frame}
\begin{frame}{Zero-width top}
\normalsize
If a staircase top exists only at \(x=4\), a grid can see it; an ideal continuous proposal hits it with probability zero. Check geometry before tuning \(p\) or \(\sigma\).
\end{frame}
\begin{frame}{Repairing the target}
\normalsize
\begin{center}\includegraphics[width=.84\textwidth,height=.53\textheight,keepaspectratio]{../../figures/mutation_reach_04_the_dead_staircase.png}\end{center}
{\small Widening the top band to 0.5 lets the same search rule reach it.}
\end{frame}
\begin{frame}{Compare moves}
\normalsize
\begin{center}\includegraphics[width=.84\textwidth,height=.53\textheight,keepaspectratio]{../../figures/permutation_mutation_02_rearrangements.png}\end{center}
{\small Changed positions and index travel rank these operators differently.}
\end{frame}
\begin{frame}{Transfer to Planta Física}
\normalsize
For 48 proportions, choose \(p\) and \(\sigma\) on the decision scale. Gene clipping does not prove annual quota, potable minimum, or monthly continuity; use the shared evaluator.
\end{frame}
\end{document}
